\numberlesssection{ЗАКЛЮЧЕНИЕ}

Цель курсовой работы достигнута: разработана базы данных для системы онлайн-заказов еды, а также приложение для доступа к ней.

При этом в ходе работы были выполнены все поставленные задачи:
\begin{itemize}
    \item проведён анализ существующих систем онлайн-заказов еды;
    \item определены требования и ограничения к разрабатываемой базе данных и приложению;
    \item спроектирована база данных и заданы ограничения целостности;
    \item выбраны средства реализации базы данных и приложения;
    \item реализованы спроектированные база данных и приложение;
    \item проведёно исследование влияния внешнего кэширования и индексации на время выполнения запросов к базе данных.
\end{itemize}

Разработанная база данных включает 13~таблиц, реализует ролевую модель разграничения доступа, содержит хранимую функцию расчёта итоговой стоимости заказа и триггеры автоматического обновления рейтингов. Проведённые исследования показали, что время выполнения запроса к PostgreSQL растёт линейно с увеличением объёма выборки, тогда как Redis обеспечивает значительно меньший прирост. Составной индекс ускоряет выполнение запросов в 2.6--5.4~раза в зависимости
от объёма таблицы.